\documentclass[12pt]{article}

\usepackage{amsmath, amssymb, amsthm}
\usepackage{enumerate}
\usepackage{hyperref}
\usepackage{geometry}
\usepackage{newtxtext}
\usepackage{newtxmath}
\geometry{margin=1in}

\title{Finite Groups - Lecture 1}
\author{}
\date{\today}

\begin{document}
\maketitle

I assume knowledge of group theory. 
This problem is based on a series of expositions on Sylow theory and group theory by Qiaochu Yuan's blog. 
I also had in mind a proof provided by \texttt{jagr2808} on Mathematics Discord using these ideas.

\section*{Questions}

\begin{enumerate}[(a)]
  \item Describe what a finite $p$-group $G$ is and give an example. 
  Describe what a group action is, 
  and give an example for a finite $p$-group $G$ acting on a finite set, say $X$.
  Describe the subset of the finite set $X$ that is fixed by the group action of 
  a finite $p$-group $G$, this is denoted as $X^G$.
  \item Let $G$ be a finite $p$-group and $X$ a finite set on which $G$ acts.
  We will slowly build up to the proof of the following result: that $|X^G| \equiv |X| \pmod {p}$
  and in particular if $p \nmid |X|$ then $X^G$ is non-empty i.e. $X^G \neq \emptyset$ and the action of $G$ on $X$
  has a fixed point. (You may prove this now or you may proceed to the next question and come back to this question later.)
  \item (Exercise for the proof of question (b)): Let $G$ be a finite $p$-group.
  Describe $X$ as a disjoint union of $G$-orbits and the form of each $G$-orbit as a quotient of $G$ by a subgroup $H$ of $G$. Compute the cardinality of each $G$-orbit and 
  use this to show that $|X| \equiv |X^G| \pmod{p}$. What does this imply about the existence of certain orbits?
  What sort of orbits do we have in this case? 
  What does this say about the fixed points of the group action of $G$ on the set $X$?
  Hence prove the result in question (b).
  \item (Cauchy's Theorem) Let $G$ be a finite group and $p$ a prime dividing $|G|$.
  Then prove that $G$ has an element of order $p$, using $(b)$ by having a finite $p$-group act on a finite set in a clever way. What set is used? Why is its cardinality of this finite set important?

  \item The exercise is to fill in the proof sketch that I saw for the Sylow theorems.
  
  (Theorem 1) Consider a $p$-subgroup $P$ of $G$, let $P$ act on $G/P$ by left multiplication. Compute the fixed points. Apply induction and Cauchy's theorem (with some extra work) on these fixed points, and prove that Sylow $p$-subgroups exist.
  \item (Theorem 3) If $P$ is a Sylow $p$-subgroup of $G$, using the conjugation action on Sylow $p$-subgroups, compute $[G:N_G(P)]$. Hence, compute the number of conjugates of $P$.
  \item (Theorem 2) Let $P$ and $P'$ be two different Sylow $p$-subgroups. Let $P$ act on $G/P'$ by left multiplication. Compute the number of fixed points and hence show that $P$ and $P'$ are conjugate.
\end{enumerate}

\newpage

\section*{Solutions}

\begin{enumerate}[(a)]

\item Let $G$ be a finite $p$-group acting on a finite set $X$.
Decompose $X$ into $G$-orbits:
\[
X=\bigsqcup_i Gx_i .
\]
For $x\in X$,
\[
Gx\cong G/G_x,
\]
where
\[
G_x=\{g\in G:gx=x\}.
\]
Hence
\[
|Gx|=[G:G_x],
\]
which is a power of $p$.

An orbit has size $1$ precisely when $G_x=G$, i.e. when $x\in X^G$.
All other orbits have size divisible by $p$. Therefore
\[
|X|=|X^G|+\sum_{\text{non-fixed orbits}}|Gx|
\]
and hence
\[
|X|\equiv |X^G|\pmod p.
\]

If $p\nmid |X|$, then $|X^G|\not\equiv0\pmod p$, so
\[
X^G\neq\emptyset .
\]

\item Let $G$ act on itself by conjugation:
\[
g\cdot x=gxg^{-1}.
\]
The fixed points are exactly the center:
\[
X^G=Z(G).
\]
By the previous result,
\[
|Z(G)|\equiv |G|\pmod p.
\]
Since $|G|$ is divisible by $p$,
\[
p\mid |Z(G)|.
\]
As $e\in Z(G)$, we have
\[
Z(G)\neq\{e\}.
\]

Now use induction on $|G|$. Since $Z(G)\neq\{e\}$,
\[
|G/Z(G)|<|G|.
\]
Also $G/Z(G)$ is a finite $p$-group, so by induction it is nilpotent.
Hence $G$ is nilpotent.

\item Let $p\mid |G|$ and define
\[
X=\{(a_1,\dots,a_p)\in G^p:
a_1a_2\cdots a_p=e\}.
\]
Then
\[
|X|=|G|^{p-1}.
\]

Let $C_p$ act on $X$ by cyclic permutation:
\[
(a_1,\dots,a_p)\mapsto(a_2,\dots,a_p,a_1).
\]
The fixed points satisfy
\[
a_1=\cdots=a_p=a
\]
and therefore
\[
a^p=e.
\]

By the fixed point theorem,
\[
|X^{C_p}|\equiv |X|\equiv0\pmod p.
\]
Since $(e,\dots,e)$ is one fixed point, there is another:
\[
(a,\dots,a)
\]
with $a\neq e$.
Thus
\[
a^p=e
\]
and the order of $a$ is $p$.

\item Let $P$ be a $p$-subgroup of $G$. Let $P$ act on
\[
G/P
\]
by left multiplication.

A coset $gP$ is fixed exactly when
\[
xgP=gP\qquad\forall x\in P,
\]
which is equivalent to
\[
g^{-1}Pg\subseteq P.
\]
Hence
\[
(G/P)^P=N_G(P)/P.
\]

By the fixed point theorem,
\[
[N_G(P):P]\equiv[G:P]\pmod p.
\]

Using induction and Cauchy's theorem on $N_G(P)/P$, one obtains a larger
$p$-subgroup unless $P$ is maximal. Therefore a maximal $p$-subgroup exists,
and it is a Sylow $p$-subgroup.

\item Let $P$ be a Sylow $p$-subgroup. The conjugates of $P$ are counted by
\[
[G:N_G(P)].
\]

Also
\[
[G:P]=[G:N_G(P)][N_G(P):P].
\]
Because $P$ is Sylow,
\[
[N_G(P):P]\equiv1\pmod p.
\]
Hence
\[
[G:N_G(P)]\equiv[G:P]\pmod p.
\]

Therefore the number of Sylow $p$-subgroups is congruent to
\[
[G:P]\pmod p.
\]

\item Let $P,P'$ be Sylow $p$-subgroups. Let $P$ act on
\[
G/P'
\]
by left multiplication.

A fixed point satisfies
\[
g^{-1}Pg\subseteq P'.
\]
Both groups are Sylow $p$-subgroups, so they have the same order, hence
\[
g^{-1}Pg=P'.
\]
Therefore
\[
P'=g^{-1}Pg,
\]
so all Sylow $p$-subgroups are conjugate.

\end{enumerate}

\end{document}